% ============================================================
% Page 3 — Alban à l'aube de notre histoire
% Transcription validée du document source.
% Mise en page simple ; typographie et lettrines à traiter
% ultérieurement dans le préambule et le modèle général.
% ============================================================

\thispagestyle{plain}

\vspace*{12mm}

\begin{center}
    {\Large\bfseries Alban à l'aube de notre histoire}

    \vspace{2mm}
    {\itshape (relais gallo-romain)}
\end{center}

\vspace{18mm}

\noindent
Alban fait d'abord partie de la province romaine (Provence).
Avec l'empereur Octave Auguste, il est compris dans l'Aquitaine.
Et après avoir été englobé dans le pays des Ruthènes, le voilà à la
fin du IVème siècle échouant dans la « Civitas Albiensium » limitée
par le Viaur, l'Aveyron, le Thoré, l'Agout.

\par\vspace{9mm}

\noindent
A cette même époque et mêmes limites naît probablement le diocèse
d'Albi. Les découvertes de l'Abbé Hermet de Rodez et de M. Sahue de
Saint Pons tout le long de la route Albi-Béziers, dalles en parties
conservées, ornières encore existantes, vestiges de villas --
statuettes -- disent assez combien les Gallo-romains fréquentaient
cette voie. Alban servit sans doute de « mansio », gîte d'étape ou
relais en même temps qu'il était le débouché naturel de toute la
région minière des environs.

\vfill
